% Centralized, decentralized and distributed control (processing units in red).
\begin{tikzpicture}[font=\sffamily\small,
    robot/.style={circle, draw=primary, fill=primary!25, minimum size=3.2mm, inner sep=0pt},
    cpu/.style={circle, draw=accent, fill=accent!80, minimum size=4.4mm, inner sep=0pt},
    link/.style={primary!60, thick}]
  % (a) centralized
  \begin{scope}
    \node[cpu] (c) at (0,0) {};
    \foreach \a in {0,45,...,315} {
      \node[robot] (r\a) at (\a:1.25) {};
      \draw[link] (c) -- (r\a);
    }
    \node at (0,-2.0) {(a) centralized};
  \end{scope}
  % (b) decentralized: several local hubs, each serving its own robots
  \begin{scope}[xshift=4.6cm]
    \foreach \h/\x/\y in {1/-0.9/0.7, 2/0.95/0.55, 3/0/-0.85} {
      \node[cpu] (h\h) at (\x,\y) {};
      \foreach \a in {0,72,...,288} {
        \node[robot] (b\h-\a) at ($(\x,\y)+(\a+20*\h:0.75)$) {};
        \draw[link] (h\h) -- (b\h-\a);
      }
    }
    \draw[link, accent!70] (h1) -- (h2) -- (h3) -- (h1);
    \node at (0,-2.0) {(b) decentralized};
  \end{scope}
  % (c) distributed: every robot is a processing unit
  \begin{scope}[xshift=9.0cm]
    \foreach \n/\x/\y in {1/-0.9/0.8, 2/0.2/1.0, 3/1.0/0.3, 4/-0.2/0.05, 5/-1.1/-0.5,
                          6/0.1/-0.9, 7/1.1/-0.8}
      \node[cpu] (d\n) at (\x,\y) {};
    \foreach \a/\b in {1/2, 2/3, 1/4, 2/4, 3/4, 1/5, 4/5, 4/6, 5/6, 6/7, 3/7, 4/7}
      \draw[link] (d\a) -- (d\b);
    \node at (0,-2.0) {(c) distributed};
  \end{scope}
\end{tikzpicture}
